\documentclass[11pt]{article}

% --------------------------------------------------
% Packages
% --------------------------------------------------
\usepackage[T1]{fontenc}
\usepackage{lmodern}
\usepackage{geometry}
\usepackage{microtype}
\usepackage{setspace}
\usepackage{amsmath,amssymb,amsthm,mathtools}
\usepackage{csquotes}
\usepackage{hyperref}
\usepackage{booktabs}

\geometry{margin=1in}
\setstretch{1.15}

% --------------------------------------------------
% Theorem Environments
% --------------------------------------------------
\newtheorem{definition}{Definition}
\newtheorem{proposition}{Proposition}
\newtheorem{remark}{Remark}

% --------------------------------------------------
% Macros
% --------------------------------------------------
\newcommand{\State}{\mathcal{S}}
\newcommand{\Event}{\mathcal{E}}
\newcommand{\Future}{\mathcal{F}}
\newcommand{\Noise}{\mathcal{N}}
\newcommand{\Graph}{\mathcal{G}}

% --------------------------------------------------
% Title
% --------------------------------------------------
\title{Event-Nodes, Texture-Nodes, and Causal Structure:\\
From State Intervention to Event-Historical Mechanistic Explanation}
\author{Flyxion}
\date{December 2025}

\begin{document}
\maketitle

% ==================================================
\section{Introduction}
% ==================================================

Recent advances in mechanistic interpretability have emphasized causal reasoning as the criterion separating genuine understanding from descriptive inspection. Interventions, counterfactuals, and causal abstractions offer a disciplined way to test whether a hypothesized algorithm is truly implemented by a system rather than merely correlated with its behavior.

However, standard causal–mechanistic approaches remain implicitly state-centric. Interventions are typically defined as instantaneous value substitutions, and irreversibility enters only through accidental information loss (e.g.\ noise or ablation). This essay argues that a minimal extension of the framework is required to account for history-sensitive structure, irreversible commitments, and the emergence of constraints that persist beyond any single state.

To this end, we introduce two complementary causal primitives:
\emph{event-nodes} and \emph{texture-nodes}. These notions allow causal graphs to represent not only values and pathways, but also the historical processes and structural textures that shape future behavior.

---

% ==================================================
\section{Causal Graphs and State-Based Intervention}
% ==================================================

Let a system be represented as a directed acyclic graph
\[
\Graph = (V, E),
\]
where vertices $V$ correspond to internal variables and edges $E$ represent causal influence. Traditional interventions replace the value of a node $v \in V$ with an externally supplied value $v'$, breaking incoming edges while preserving outgoing ones.

This formalism excels at identifying mediating variables and testing necessity or sufficiency within a fixed causal structure. However, it treats causal nodes as temporally atomic: values are swapped, tested, and restored without regard to how those values were produced.

Such a framework can identify \emph{paths of influence}, but not \emph{paths of becoming}.

---

% ==================================================
\section{Event-Nodes}
% ==================================================

\begin{definition}[Event-Node]
An \emph{event-node} is a causal node whose primary explanatory role is not the value it holds, but the irreversible constraint it imposes on the system's future admissible states.
\end{definition}

Formally, an event-node $e \in \Event$ induces a transition
\[
\Future_{t+1} = \Future_t \cap \Phi_e,
\]
where $\Phi_e \subseteq \State$ is a constraint set defining which future states remain reachable.

Unlike state variables, event-nodes cannot be fully characterized by their instantaneous values. Their causal significance lies in the \emph{reduction of future possibility space}. An intervention on an event-node is therefore not merely a substitution, but a commitment.

This allows irreversibility to be treated as structural rather than accidental. Even if all information is preserved, an event-node may permanently foreclose certain counterfactual trajectories.

---

% ==================================================
\section{Texture-Nodes}
% ==================================================

\begin{definition}[Texture-Node]
A \emph{texture-node} is a causal node that parametrizes the generative structure of variation within a system, shaping how local perturbations propagate without directly determining outcomes.
\end{definition}

Texture-nodes do not encode discrete decisions or constraints. Instead, they modulate \emph{how} behavior is generated. Their role is analogous to material or shader parameters in a three-dimensional mesh: they do not alter topology, but they determine surface character, continuity, and response to deformation.

Let $\Noise : \State \times \Theta \to \State$ be a noise or variation generator. A texture-node corresponds to a parameter vector $\theta \in \Theta$ such that
\[
s_{t+1} = \Noise(s_t, \theta).
\]

Crucially, $\theta$ may itself be the result of prior events. Thus texture-nodes often record ratios, gradients, or directional biases accumulated over time.

---

% ==================================================
\section{Analogy: Meshes, Materials, and Generative Noise}
% ==================================================

Consider a geometric mesh in a modeling environment. The mesh defines topology; materials define appearance. Changing a material does not move vertices, yet it radically alters how light, texture, and noise interact with the surface.

In this analogy:
\begin{itemize}
\item State-nodes correspond to vertex positions.
\item Event-nodes correspond to irreversible topology edits.
\item Texture-nodes correspond to material parameters.
\item Noise generators correspond to procedural variation functions.
\end{itemize}

A texture-node selects a generator function whose output depends on a ratio or vector parameter:
\[
\theta = \frac{\alpha}{\beta},
\]
where $\alpha$ and $\beta$ may represent accumulated gradients, pressures, or historical frequencies. The system records $\theta$ not as a symbolic label, but as a bias shaping future dynamics.

Thus, two systems with identical topology and instantaneous states may diverge due to differing textures: the same perturbation produces different outcomes because the generative field differs.

---

% ==================================================
\section{Counterfactuals Over Futures}
% ==================================================

Traditional counterfactuals ask:
\begin{quote}
What would the output have been if this variable had taken a different value?
\end{quote}

Event- and texture-nodes enable a stronger question:
\begin{quote}
What futures would still be reachable if this event or texture had been different?
\end{quote}

Let $\mathcal{R}(s)$ denote the reachable future set from state $s$. An event-level counterfactual compares $\mathcal{R}(s \mid e)$ to $\mathcal{R}(s \mid e')$. A texture-level counterfactual compares the shape of $\mathcal{R}(s)$ under different generative parameters $\theta$.

This reframes semantic compatibility: a substitution fails not because a value is wrong, but because it presupposes a future space that no longer exists.

---

% ==================================================
\section{Causal Abstraction Revisited}
% ==================================================

A causal abstraction is traditionally validated by interventional equivalence. With event- and texture-nodes, abstraction gains an additional requirement:

\begin{definition}[Temporal Coherence]
A causal abstraction is temporally coherent if it preserves both interventional outcomes and the irreversibility structure of the underlying system.
\end{definition}

An abstraction that predicts reversibility where the system exhibits irreversible settlement is incomplete, even if its input–output behavior matches locally.

This explains why multiple abstractions may coexist: each compresses structure differently, preserving certain textures while discarding others.

---

% ==================================================
\section{Failure, Collapse, and Mixture}
% ==================================================

Failures under intervention are not noise; they are diagnostic. A collapse into incoherence indicates that a texture-node or event-node essential to generative stability has been disrupted.

Systems that exhibit patchwork behavior under intervention are best understood as mixtures of event-histories and textures rather than single algorithms. Their apparent inconsistency reflects genuine structural plurality.

---

% ==================================================
\section{What It Means to Understand}
% ==================================================

Understanding is not achieved by control alone, nor by prediction in isolation. It consists in reconstructing:
\begin{itemize}
\item the event-nodes that constrain futures,
\item the texture-nodes that shape variation,
\item and the causal graph that organizes both.
\end{itemize}

Such understanding is necessarily partial. Texture is lossy under abstraction, and event-history cannot be fully replayed. But these limitations are principled, not failures.

---

% ==================================================
\section{Conclusion}
% ==================================================

By extending causal graphs with event-nodes and texture-nodes, mechanistic interpretability gains the expressive power to describe irreversible structure, generative variation, and historical constraint without abandoning causal rigor.

The result is a framework in which mechanisms are not merely paths through a graph, but histories inscribed into the future—where what a system \emph{can no longer do} is as explanatory as what it does.

% ==================================================
\section{Necessary Conditions for Texture-Nodes to Become Event-Nodes}
% ==================================================

This section derives conditions under which a variable that initially functions
as a \emph{texture-node}---a parameter modulating a generative field without
directly constraining the admissible futures---must be reclassified as an
\emph{event-node}, i.e.\ as a node whose causal significance is the irreversible
restriction of future reachable states.

\subsection{Preliminaries: Reachability and Generative Fields}

Let $\State$ be the system state space and let $\mathcal{R}_T(s)$ denote the
set of states reachable from $s \in \State$ within a horizon $T \in \mathbb{N}$.
We assume a one-step update law of the form
\begin{equation}
s_{t+1} = F(s_t;\theta_t),
\end{equation}
where $\theta_t \in \Theta$ is a parameter vector. When $\theta_t$ is a
texture-node, its intended role is to modulate the generative field (e.g.\ the
noise, smoothness, anisotropy, sparsity, or routing biases) without imposing
hard, history-settling constraints.

To make the modulation explicit, we may decompose
\begin{equation}
F(s;\theta) = G(s) + H(s;\theta),
\end{equation}
where $G$ is the ``baseline'' dynamics and $H$ is the texture-controlled
component (which may be stochastic or deterministic).

\begin{definition}[Texture-Equivalence of Futures]
Two parameter values $\theta,\theta' \in \Theta$ are \emph{texture-equivalent at
horizon $T$} for state $s$ if
\begin{equation}
\mathcal{R}_T(s \mid \theta) = \mathcal{R}_T(s \mid \theta').
\end{equation}
They are \emph{approximately texture-equivalent} if the reachable sets are
equal up to a specified tolerance (e.g.\ Hausdorff distance or probability mass
distance under a stochastic transition kernel).
\end{definition}

Under pure texture modulation, changes in $\theta$ may alter the \emph{measure}
over futures (which trajectories are likely), while preserving the \emph{support}
of futures (which trajectories are possible). The distinction between measure
and support is the fulcrum of the event/texture boundary.

\subsection{Support Restriction Criterion}

Assume the system admits a stochastic kernel representation
\begin{equation}
s_{t+1} \sim K_{\theta_t}(\,\cdot \mid s_t),
\end{equation}
so that for any measurable set $A \subseteq \State$,
\begin{equation}
\mathbb{P}(s_{t+1} \in A \mid s_t) = \int_A K_{\theta_t}(ds' \mid s_t).
\end{equation}

\begin{definition}[Support of the One-Step Kernel]
The \emph{support} of $K_{\theta}(\cdot \mid s)$ is the smallest closed set
$\mathrm{supp}\,K_{\theta}(\cdot \mid s) \subseteq \State$ with
probability one.
\end{definition}

\begin{proposition}[Necessary Condition: Support Change]
If a texture-node value $\theta$ is to remain a pure texture-node (i.e.\ not an
event-node), then for all relevant states $s$ and horizons $T$,
changes in $\theta$ may alter likelihoods but must not remove futures from the
reachable support. In particular, it is necessary that
\begin{equation}
\mathrm{supp}\,K_{\theta}(\cdot \mid s) = \mathrm{supp}\,K_{\theta'}(\cdot \mid s)
\quad \text{for all } s \text{ in the regime of interest.}
\end{equation}
\end{proposition}

\begin{proof}
If $\mathrm{supp}\,K_{\theta}(\cdot \mid s) \neq \mathrm{supp}\,K_{\theta'}(\cdot \mid s)$,
then there exists a set $A$ such that one parameter value assigns nonzero
probability to transitioning into $A$ while the other assigns zero probability.
This implies a strict inclusion between one-step reachable sets, and therefore
a strict inclusion between $T$-step reachable sets for some $T$ (at worst $T=1$).
Hence the parameter change removes admissible futures and functions as a
constraint on reachability, which is the defining role of an event-node.
\end{proof}

Intuitively, a texture parameter may act like a material selection that changes
surface appearance, but if it changes the mesh's admissible deformations (e.g.\
turns a soft body into a rigid body that cannot bend), then it has become a
structural operator and thus an event.

\subsection{Irreversibility Criterion via Hysteresis}

Even if a parameter does not change one-step support, it may still be
event-like if its \emph{update law} introduces hysteresis or path dependence.

Assume the texture parameter evolves according to
\begin{equation}
\theta_{t+1} = U(\theta_t, s_t),
\end{equation}
where $U$ may include thresholding, clipping, quantization, or saturation.

\begin{definition}[Hysteretic Texture Update]
The update $U$ is \emph{hysteretic} on a set $\Omega \subseteq \Theta \times \State$
if there exist two histories $(\theta_0,s_0,\dots,s_t)$ and
$(\theta_0',s_0',\dots,s_t')$ such that $(\theta_t,s_t)=(\theta_t',s_t')$ but
$\theta_{t+1} \neq \theta_{t+1}'$.
\end{definition}

\begin{proposition}[Necessary Condition: No Hysteresis Under State Equivalence]
If $\theta$ is to remain a pure texture-node, it is necessary that its update
law be \emph{non-hysteretic} under the state equivalences used in causal
analysis. Concretely, whenever the causal framework treats two situations as
equivalent (e.g.\ identical activations at the intervention boundary), the
texture update must agree:
\begin{equation}
(\theta_t,s_t) \sim (\theta_t',s_t') \implies U(\theta_t,s_t)=U(\theta_t',s_t').
\end{equation}
\end{proposition}

\begin{proof}
If two situations are treated as equivalent by the causal testing protocol but
yield different $\theta_{t+1}$, then the future reachable sets diverge purely
due to untracked historical information. In that case $\theta$ encodes
irreducible history, and interventions that only set instantaneous values are
insufficient to restore counterfactual symmetry. The node is therefore
functioning as a history-settling operator, i.e.\ an event-node.
\end{proof}

Hysteresis corresponds to a material parameter whose value depends on
\emph{how} it was reached, not merely its current numeric value. In procedural
shading terms, this is the difference between a parameter that selects a noise
seed and a parameter that has been permanently baked into a texture map.

\subsection{Thresholding and Phase Boundary Condition}

Many systems contain sharp boundaries where a continuously varying texture
parameter induces a qualitative shift in reachable structure.

Suppose there exists a scalar summary $\rho(\theta)$ (for example, a ratio of
two accumulated quantities) such that the dynamics effectively changes regime
when $\rho(\theta)$ crosses a threshold $\rho_\ast$:
\begin{equation}
F(s;\theta) =
\begin{cases}
F_-(s;\theta), & \rho(\theta) < \rho_\ast,\\
F_+(s;\theta), & \rho(\theta) \ge \rho_\ast.
\end{cases}
\end{equation}

\begin{proposition}[Necessary Condition: No Regime-Switching Under Small Perturbations]
For $\theta$ to remain texture-like on a regime of interest, it is necessary
that small perturbations in $\theta$ not induce regime switching that changes
reachable support. In particular, for all $(s,\theta)$ in the regime and all
sufficiently small $\delta$,
\begin{equation}
\mathcal{R}_T(s \mid \theta) \approx \mathcal{R}_T(s \mid \theta+\delta)
\end{equation}
in support, not merely in probability.
\end{proposition}

\begin{proof}
If arbitrarily small perturbations can flip the system between distinct
transition regimes whose reachable supports differ, then the parameter acts as
a gate selecting among different future constraint families. Such gating is
structural and thus event-like.
\end{proof}

This is the minimal formal statement of a ``phase boundary.'' Once a texture
parameter gates between distinct dynamical laws, it ceases to be mere texture.

\subsection{Recording Criterion: When a Parameter Becomes a Commitment}

Your guiding analogy includes that the texture generator depends on a recorded
ratio that functions as a vector parameter for a variable in question. Let
\begin{equation}
\theta_t = \Psi\!\left(\frac{a_t}{b_t}\right),
\end{equation}
where $a_t,b_t$ are accumulated quantities (counts, energies, gradients,
frequencies) and $\Psi$ maps the ratio into a vector parameter (direction,
scale, anisotropy, routing preference).

If $a_t$ and $b_t$ are merely summaries that can be recomputed from current
state, then $\theta_t$ remains texture-like. If instead $a_t$ and $b_t$ encode
irrecoverable history, then $\theta_t$ inherits event-like character.

\begin{proposition}[Necessary Condition: Reconstructibility of the Recorded Ratio]
If $\theta_t$ is computed from a recorded ratio $\frac{a_t}{b_t}$, then a
necessary condition for $\theta$ to remain a texture-node is that the ratio be
reconstructible from the instantaneous causal boundary used in interventions.
That is, there must exist a function $\Xi$ such that
\begin{equation}
\frac{a_t}{b_t} = \Xi(s_t)
\end{equation}
on the regime of interest (or, in the stochastic case, that its conditional
distribution depends only on $s_t$).
\end{proposition}

\begin{proof}
If the ratio cannot be reconstructed from the intervention boundary, then two
states treated as identical under the causal protocol may still carry
different latent ratios, hence different $\theta_t$, hence different future
supports or constraints. The parameter then encodes untracked history that
cannot be erased by state-level substitution, making it a commitment variable,
i.e.\ event-like.
\end{proof}

In the mesh analogy, reconstructibility corresponds to whether the material
choice is a pure function of current shader inputs, versus a baked selection
made earlier in the pipeline that cannot be inferred from the final rendered
surface alone.

\subsection{Summary: Minimal Necessary Conditions}

The boundary between texture-nodes and event-nodes is not metaphysical. It is
determined by reachability, reconstructibility, and irreversibility. A node
initially treated as a texture-node must be promoted to an event-node whenever
any of the following necessary conditions fail in the regime of interest:

\begin{enumerate}
\item \textbf{Support invariance:} changing the parameter does not change the
support of the transition kernel (no futures become impossible).
\item \textbf{No hysteresis:} the parameter update does not depend on untracked
history under the equivalences used by intervention tests.
\item \textbf{No regime gating:} small changes in the parameter do not switch
the system between distinct dynamical regimes with different reachable supports.
\item \textbf{Reconstructible recording:} any recorded ratio or summary that
drives the parameter can be reconstructed from the causal boundary state (or
its conditional distribution depends only on that boundary state).
\end{enumerate}

When these conditions hold, the parameter behaves as a genuine texture: it
reshapes variation without settling futures. When any condition fails, the
parameter is no longer merely a material choice; it becomes a structural edit
to the system's admissible counterfactual space, and therefore functions as an
event-node.

% ==================================================
\section{Sufficiency Conditions and the Event--Texture Phase Boundary}
% ==================================================

The previous section derived \emph{necessary} conditions under which a
texture-node must be promoted to an event-node. We now derive complementary
\emph{sufficiency} conditions: conditions under which a variable can be treated
as a genuine texture-node without loss of causal adequacy within a given
regime. Together, these results define a sharp phase boundary between texture
and event structure.

\subsection{Texture Sufficiency via Future-Space Equivalence}

Let $\mathcal{R}_\infty(s \mid \theta)$ denote the infinite-horizon reachable
future set from state $s$ under parameter $\theta$. While practical analysis
often works at finite horizons, the conceptual distinction between texture and
event is asymptotic: event-nodes settle futures permanently.

\begin{definition}[Future-Space Equivalence Class]
Two parameter values $\theta,\theta' \in \Theta$ are said to lie in the same
\emph{future-space equivalence class} on a regime $\Omega \subseteq \State$ if
\begin{equation}
\mathcal{R}_\infty(s \mid \theta) = \mathcal{R}_\infty(s \mid \theta')
\quad \text{for all } s \in \Omega.
\end{equation}
\end{definition}

\begin{proposition}[Sufficiency Condition: Asymptotic Equivalence]
If a parameter $\theta$ varies only within a single future-space equivalence
class on the regime of interest, then $\theta$ can be treated as a texture-node
for all causal abstractions restricted to that regime.
\end{proposition}

\begin{proof}
By hypothesis, variations in $\theta$ do not alter the set of admissible
infinite-horizon futures. Consequently, no irreversible commitments are
introduced by changing $\theta$, and all counterfactual trajectories remain
reachable under replay. Thus $\theta$ does not impose structural constraints
and functions purely as a modulator of trajectory selection within a fixed
future space, which is precisely the role of a texture-node.
\end{proof}

This formalizes the intuitive notion that texture governs \emph{how} futures are
sampled, not \emph{which} futures exist.

\subsection{Local Reversibility and Retexturing}

A second sufficiency condition concerns reversibility under intervention.

\begin{definition}[Local Retexturability]
A parameter $\theta$ is \emph{locally retexturable} at state $s$ if, for any
$\theta'$ in the parameter domain, there exists a finite sequence of
state-level interventions $\{I_k\}$ such that
\begin{equation}
(s,\theta) \xrightarrow{I_1,\dots,I_n} (s',\theta')
\end{equation}
with $\mathcal{R}_\infty(s') = \mathcal{R}_\infty(s)$.
\end{definition}

\begin{proposition}[Sufficiency Condition: Retexturability]
If a parameter is locally retexturable throughout the regime of interest, then
it cannot function as an event-node within that regime.
\end{proposition}

\begin{proof}
Local retexturability guarantees that any value of $\theta$ can be imposed
without altering the infinite-horizon future space. Thus no parameter setting
constitutes a point of no return. Since event-nodes are defined by irreversible
constraint of futures, a locally retexturable parameter lacks the defining
property of an event-node.
\end{proof}

In the mesh analogy, this corresponds to a material parameter that can always be
swapped at render time, never requiring a topology rebuild or geometry bake.

\subsection{Texture Stability Under Causal Abstraction}

A further sufficiency condition concerns abstraction itself.

\begin{definition}[Abstraction-Stable Texture]
A parameter $\theta$ is \emph{abstraction-stable} if, for every valid causal
abstraction $\mathcal{A}$ of the system on the regime of interest, the image of
$\theta$ under $\mathcal{A}$ does not correspond to an event-node in the
abstracted graph.
\end{definition}

\begin{proposition}[Sufficiency Condition: Abstraction Stability]
If a parameter is abstraction-stable across all abstractions admissible for the
intended explanatory task, then treating it as a texture-node is sufficient for
mechanistic understanding at that level.
\end{proposition}

\begin{proof}
If no abstraction relevant to the explanatory task requires $\theta$ to be
represented as an irreversible constraint, then $\theta$ plays no essential
role in determining the abstract system's commitment structure. Representing
$\theta$ as a texture-node preserves all explanatory distinctions made at that
level, and promoting it to an event-node would add structure without
explanatory gain.
\end{proof}

This result makes explicit that texture versus event is not an absolute
classification, but a level-relative one.

\subsection{Phase Boundary Summary}

Combining necessary and sufficiency conditions, we obtain a precise phase
boundary:

\begin{remark}[Event--Texture Phase Boundary]
A causal variable functions as a texture-node on a regime of interest if and
only if:
\begin{enumerate}
\item its variation preserves the support of reachable futures,
\item it exhibits no hysteresis under the equivalence relations used by
intervention protocols,
\item it does not gate regime changes under small perturbations,
\item any recorded summaries driving it are reconstructible at the causal
boundary,
\item its values lie within a single future-space equivalence class,
\item and it is locally retexturable without altering future reachability.
\end{enumerate}
Violation of any of these conditions forces reclassification as an event-node
for that regime.
\end{remark}

The distinction is thus neither metaphysical nor merely descriptive. It is a
formal consequence of how a system partitions its future. Texture governs
variation within a future; event governs which futures remain possible.

This phase boundary provides a principled basis for extending causal
mechanistic analysis beyond instantaneous state substitution to systems whose
behavior is shaped by accumulated structure, recorded ratios, and irreversible
settlement.

% ==================================================
\section{Connection to Spherepop Calculus}
% ==================================================

The distinction between event-nodes and texture-nodes admits a direct and
natural correspondence with the primitives of Spherepop calculus. This section
makes that correspondence explicit and shows that the extended causal framework
developed above can be seen as a semantic realization of Spherepop’s operational
rules.

\subsection{Event-Nodes as \textsc{Pop} Operations}

In Spherepop calculus, a \textsc{Pop} is defined as an irreversible operation
that eliminates a subset of admissible futures. This is not a change of state
in the extensional sense, but a settlement of possibility.

Event-nodes coincide exactly with this notion.

Formally, let $\mathcal{H}_t$ denote the admissible future set at time $t$. A
Spherepop \textsc{Pop} induces a transition
\begin{equation}
\mathcal{H}_{t+1} = \mathcal{H}_t \setminus \Delta,
\end{equation}
for some nonempty exclusion set $\Delta$.

This is identical to the action of an event-node as defined earlier:
\begin{equation}
\Future_{t+1} = \Future_t \cap \Phi_e,
\end{equation}
with $\Phi_e = \State \setminus \Delta$.

Thus:
\begin{center}
\textbf{Event-node} $\;\longleftrightarrow\;$ \textbf{Spherepop \textsc{Pop}}
\end{center}

Both are characterized by:
\begin{itemize}
\item irreversibility,
\item future-space restriction,
\item and independence from instantaneous representational value.
\end{itemize}

Crucially, neither requires destructive noise or loss of information. The
irreversibility is structural, not epistemic.

\subsection{Texture-Nodes as Spherepop \textsc{Bind} Structure}

Spherepop calculus distinguishes \textsc{Bind} operations from \textsc{Pop}s.
A \textsc{Bind} introduces ordering, dependency, or constraint relations
between future actions without eliminating options outright.

Texture-nodes implement exactly this role at the causal–mechanistic level.

A texture-node:
\begin{itemize}
\item reshapes the generative field,
\item biases which futures are likely,
\item and introduces anisotropy or smoothness,
\end{itemize}
while preserving the support of admissible futures.

This matches the operational role of \textsc{Bind}:
\begin{equation}
\mathcal{H}_{t+1} = \mathcal{H}_t,
\qquad
\mu_{t+1} \neq \mu_t,
\end{equation}
where $\mu_t$ is a measure or ordering over $\mathcal{H}_t$.

Thus:
\begin{center}
\textbf{Texture-node} $\;\longleftrightarrow\;$ \textbf{Spherepop \textsc{Bind}}
\end{center}

The noise-generator analogy fits precisely here. A texture-node selects or
modulates a generator function whose parameters are ratios or vectors recorded
from prior activity. These ratios do not settle futures, but they bias how
future trajectories are drawn.

\subsection{Recorded Ratios and Spherepop History}

In Spherepop calculus, objects and relations exist only insofar as they are
introduced by events. Identity is historical, not extensional.

The recorded ratios driving texture-nodes correspond to \emph{derived
historical summaries} in Spherepop: quantities accumulated from prior pops and
binds that shape subsequent dynamics without constituting new pops themselves.

Let $r_t = a_t / b_t$ be a recorded ratio influencing a texture-node parameter
$\theta_t = \Psi(r_t)$. In Spherepop terms:
\begin{itemize}
\item $a_t$ and $b_t$ are functions of the event history,
\item $r_t$ is a historical invariant,
\item $\theta_t$ is a derived binding parameter.
\end{itemize}

When $r_t$ is reconstructible from the current causal boundary, the binding
remains soft and replayable. When it is not, the ratio encodes untracked
history, and the texture-node must be promoted to an event-node. In Spherepop
terms, a \textsc{Bind} has hardened into a \textsc{Pop}.

This precisely mirrors the necessary-condition results derived earlier.

\subsection{Collapse and Spherepop \textsc{Collapse}}

Spherepop calculus defines \textsc{Collapse} as the erasure of historical
differentiation in order to restore flexibility. This operation discards event
records while preserving optionality.

In the causal framework, collapse corresponds to:
\begin{itemize}
\item resetting or marginalizing event-nodes,
\item erasing recorded ratios,
\item restoring symmetry between counterfactual futures.
\end{itemize}

Importantly, collapse is not the inverse of a pop. It does not recover lost
futures; it removes the record that distinguished them. This aligns with the
observation that even after collapse, some futures may remain inaccessible if
the underlying structure has been irreversibly altered.

\subsection{Causal Abstraction as Spherepop Replay}

Spherepop calculus emphasizes replay: the idea that meaning and identity arise
from the reproducibility of event histories.

Causal abstraction in the extended mechanistic framework can be understood as
\emph{replay equivalence}:
\begin{quote}
Two systems are abstractly equivalent if replaying corresponding sequences of
pops and binds yields the same admissible future structure.
\end{quote}

Interchange interventions correspond to replaying alternative event sequences
under controlled substitutions. Abstraction validity requires that replay at
the abstract level preserves both:
\begin{itemize}
\item input–output counterfactuals, and
\item the structure of irreversible settlement.
\end{itemize}

This reframes causal abstraction not as pattern matching between graphs, but as
compatibility between event algebras.

\subsection{Summary}

The correspondence can be summarized succinctly:
\begin{center}
\begin{tabular}{ll}
\toprule
\textbf{Causal Framework} & \textbf{Spherepop Calculus} \\
\midrule
Event-node & \textsc{Pop} \\
Texture-node & \textsc{Bind} \\
Support restriction & Loss of admissible futures \\
Recorded ratio & Historical invariant \\
Retexturing & Soft binding \\
Irreversible promotion & Bind $\rightarrow$ Pop \\
Collapse & \textsc{Collapse} \\
Causal abstraction & Replay equivalence \\
\bottomrule
\end{tabular}
\end{center}

Viewed this way, the extended causal–mechanistic framework is not merely
compatible with Spherepop calculus; it is a concrete instantiation of its
event-historical semantics. Event-nodes and texture-nodes provide the missing
bridge between abstract operational rules and empirically testable causal
interventions, grounding Spherepop’s notions of irreversibility, commitment,
and replay in mechanistic analysis.

% ==================================================
\section{Categorical and Sheaf-Theoretic Interpretation}
% ==================================================

The distinction between event-nodes and texture-nodes admits a precise
reformulation in categorical and sheaf-theoretic terms. Doing so clarifies both
the limits of state-based causal analysis and the sense in which event-historical
structure constitutes genuine additional information rather than mere
annotation. In this section, event-nodes are interpreted as categorical
operations that alter the base category of admissible futures, while
texture-nodes are interpreted as internal parameterizations living in the
fibers of a sheaf over that base.

\subsection{The Category of Admissible Futures}

Let $\mathcal{C}_t$ denote a category whose objects are admissible future
continuations from time $t$, and whose morphisms represent valid transitions
between such continuations. Composition corresponds to temporal concatenation.
At any given moment, mechanistic analysis typically assumes a fixed
$\mathcal{C}_t$, implicitly treating the space of possible futures as stable
under intervention.

An event-node, however, does not merely select an object or morphism within
$\mathcal{C}_t$. Instead, it induces a functor
\[
E : \mathcal{C}_t \longrightarrow \mathcal{C}_{t+1}
\]
that is not essentially surjective. Objects in $\mathcal{C}_t$ that are not in
the image of $E$ correspond to futures that have been irrevocably excluded.
Irreversibility, in this formulation, is categorical non-invertibility: there
exists no functor $E^{-1}$ whose composition with $E$ restores the original
category up to equivalence.

This makes precise the sense in which event-nodes are not state updates but
changes in the ambient space of possibility. They act on the \emph{indexing
category} of futures itself.

\subsection{Texture-Nodes as Fiberwise Structure}

Texture-nodes, by contrast, do not alter the base category $\mathcal{C}_t$. They
live in additional structure carried over it. To formalize this, consider a
sheaf $\mathcal{F}$ over $\mathcal{C}_t$, where to each object (future) $x$ we
associate a space $\mathcal{F}(x)$ of admissible generative parameters. Sections
of this sheaf correspond to coherent assignments of texture across compatible
futures.

A texture-node corresponds to the choice of a section (or local section) of
$\mathcal{F}$. Changing a texture-node alters which section is selected, but
does not alter the underlying objects of $\mathcal{C}_t$. In other words, the
support of futures remains fixed; only their internal geometry is reshaped.

This sheaf-theoretic perspective makes clear why texture-nodes modulate
likelihoods, smoothness, or anisotropy without settling outcomes. They change
the internal structure of fibers while preserving the base.

\subsection{Recorded Ratios as Sections with Memory}

The recorded ratios discussed earlier acquire a natural interpretation in this
setting. A ratio such as $r_t = a_t / b_t$ determines a parameter
$\theta_t = \Psi(r_t)$, which in turn selects a local section of $\mathcal{F}$.
If the ratio is reconstructible from the current causal boundary, then the
section depends only on local data and defines a genuine sheaf section.

If, however, the ratio depends on global historical information that cannot be
recovered from the boundary state, then the would-be section fails the sheaf
gluing condition. It is no longer determined by compatible local data, but by a
choice that depends on how those locals were assembled. In categorical terms,
the assignment ceases to be natural.

This failure of naturality corresponds exactly to the promotion of a
texture-node to an event-node. The parameter can no longer be treated as living
in the fibers of a fixed base category; instead, it must be reified as an
operation on the base itself.

\subsection{Event Promotion as Base Change}

The transition from texture-node to event-node can therefore be understood as a
change of base. What was previously modeled as a sheaf
$\mathcal{F} \to \mathcal{C}_t$ must now be modeled as a different base category
$\mathcal{C}_{t+1}$ with its own sheaf of textures. The original sheaf does not
extend across the event, because the objects it was defined over no longer
exist.

This clarifies why no amount of state-level intervention can undo certain
effects: once the base category has changed, interventions that assume the old
base are ill-typed. Counterfactual substitutions fail not because values are
incorrect, but because they presuppose morphisms that are no longer present.

\subsection{Causal Abstraction as Functorial Projection}

Causal abstraction can now be reformulated as the search for a functor
\[
A : \mathcal{C}_t \longrightarrow \mathcal{D}
\]
from the category of admissible futures to a simpler category $\mathcal{D}$
representing an abstract algorithm. Validity of abstraction requires that $A$
commute, up to natural isomorphism, with the event functors that define the
system’s irreversible structure.

In particular, if an abstraction collapses distinctions that are preserved by
event functors—predicting reversibility where the concrete system exhibits
categorical collapse—then it fails to be a faithful representation of the
system’s temporal structure. This provides a categorical restatement of the
temporal coherence criterion introduced earlier.

\subsection{Spherepop Reinterpreted Categorically}

Within this formalism, Spherepop calculus emerges as an algebra of base changes
and sheaf reconfigurations. \textsc{Pop} operations correspond to non-invertible
endofunctors on the category of futures. \textsc{Bind} operations correspond to
changes of section within a fixed sheaf. \textsc{Collapse} corresponds to a
quotienting or truncation operation that identifies histories while preserving
the base category.

Replay, finally, corresponds to functorial equivalence: two histories are
meaningfully the same if they induce equivalent categories of admissible
futures together with equivalent sheaves of texture.

\subsection{Interpretive Payoff}

The categorical and sheaf-theoretic formulation sharpens the central thesis of
this work. Mechanistic understanding is not merely about identifying variables
or paths, but about determining which aspects of a system live in the fibers and
which act on the base. Texture is fiberwise and replayable; event structure is
base-altering and irreversible. Confusing the two leads either to spurious
claims of control or to the false expectation that all mechanisms can be undone
by sufficiently clever intervention.

By making this distinction explicit, the framework provides a principled bridge
between causal experimentation, event-historical semantics, and abstract
mathematical structure, without inflating the ontology beyond what mechanistic
analysis strictly requires.

% ==================================================
\section{Final Synthesis: Mechanism, History, and Texture}
% ==================================================

The aim of this essay has been to reconcile two perspectives that are often
treated as incompatible: the causal–mechanistic program, with its emphasis on
intervention and counterfactual control, and an event-historical view in which
irreversibility, commitment, and accumulated structure play an essential
explanatory role. The introduction of event-nodes and texture-nodes provides the
minimal conceptual vocabulary required to unify these perspectives without
weakening either.

At the level of causal experimentation, state-based interventions remain
indispensable. They allow one to isolate mediators, trace information flow, and
test hypotheses about algorithmic structure. However, taken alone, such
interventions implicitly assume that the space of admissible futures is fixed
and replayable. This assumption holds only within regimes where all relevant
structure is texture-like—where variation reshapes likelihoods but does not
settle possibilities.

Event-nodes mark the precise point at which this assumption fails. They are not
distinguished by dramatic representational change, nor by loss of information,
but by a categorical alteration of what futures remain reachable. Their effects
cannot be undone by clever substitution because they operate at the level of
constraints rather than values. In this sense, event-nodes formalize commitment:
they explain why certain counterfactuals are no longer meaningful, not merely
why they are unlikely.

Texture-nodes occupy the complementary role. They encode how variation is
generated, how noise is structured, and how tendencies are biased without
foreclosing alternatives. The analogy to material parameters in a generative
mesh is exact: texture changes appearance, smoothness, and response, but leaves
the underlying topology intact. When texture parameters are derived from
recorded ratios or accumulated statistics, they carry history forward in a soft
and replayable form. Only when such recording becomes non-reconstructible or
gates regime change does texture harden into event.

The necessary and sufficiency conditions derived earlier make this boundary
precise. A variable must be promoted to an event-node exactly when it restricts
the support of reachable futures, introduces hysteresis under the equivalences
assumed by intervention, or encodes history that cannot be erased by state-level
replay. Conversely, variables that preserve future-space equivalence, admit
retexturing, and remain fiberwise across valid abstractions may be safely treated
as texture-nodes within the relevant regime.

The categorical and sheaf-theoretic reformulation clarifies why this distinction
is not a matter of modeling preference. Event-nodes act on the base category of
admissible futures, while texture-nodes live in sheaves over that base. Promotion
from texture to event corresponds to a change of base, not merely a refinement
of parameters. Abstractions that ignore this distinction may succeed locally yet
fail to respect the temporal coherence of the system they purport to explain.

Seen through this lens, Spherepop calculus is not an external metaphor but an
operational algebra for these distinctions. \textsc{Pop} corresponds to the
introduction of event-nodes, \textsc{Bind} to the modulation of texture, and
\textsc{Collapse} to the controlled erasure of historical differentiation. Causal
abstraction becomes replay equivalence: two systems are equivalent insofar as
their sequences of events and textures induce equivalent categories of futures.

The resulting picture of mechanistic understanding is neither purely
state-based nor irreducibly historical. Understanding consists in determining
which aspects of a system are replayable textures and which are irrevocable
events, and in constructing abstractions that preserve this distinction. What a
system can no longer do is as much a part of its mechanism as what it does.

This synthesis suggests a natural stopping point for mechanistic explanation.
Beyond this boundary, further detail does not yield deeper understanding but
merely finer texture. Within it, causal intervention, event history, and
mathematical structure align to provide a coherent account of how systems
compute, commit, and constrain their own futures.

% ==================================================
\appendix
\section{Appendix: Implementing Spherepop Operators in Forth}
% ==================================================

This appendix sketches a concrete implementation strategy for Spherepop
operators in a small, inspectable Forth system. The design goal is not to build
a complete runtime, but to demonstrate that Spherepop’s event-historical
semantics can be realized as a minimal stack language where every operation
either (i) records an irreversible event, (ii) introduces replayable structural
binding, or (iii) collapses historical differentiation. Forth is particularly
well-suited because it makes evaluation order explicit, keeps state small and
visible, and encourages dictionary-driven extension.

\subsection{Representation Strategy}

A Spherepop system is most naturally represented as two coupled structures. The
first is a persistent \emph{event log} whose entries are the primitive operations
(\textsc{Pop}, \textsc{Bind}, \textsc{Merge}, \textsc{Link}, \textsc{Collapse}).
The second is a derived \emph{present state} computed from replaying some prefix
(or a compressed summary) of the log. In practice, a Forth implementation keeps
both: the event log provides semantic ground truth, while the derived state
provides efficient execution.

We assume a simple model in which objects are integer identifiers. Each object
has an associated parent pointer or part-of relation, plus optional metadata.
The semantics of part-of can be built as a union-find variant (for merge-like
operations) supplemented by additional relations (for link-like operations).
The essential Spherepop requirement is that these relations are not treated as
axiomatic; they exist only insofar as they are generated by recorded events.

\subsection{Forth Data Structures}

We implement a small heap using arrays. The details can be varied across Forth
systems; the following style is compatible with common ANS Forth conventions.

We represent:
\begin{enumerate}
\item a monotonically increasing object counter,
\item an event log with fixed-width records,
\item parent/representative arrays for merge semantics,
\item an optional adjacency list region for link semantics.
\end{enumerate}

The event log records a tag and up to three integer fields. The tag determines
interpretation. This deliberately mirrors the idea that the log is authoritative,
and the present state is derived.

\begin{verbatim}
\ ----------------------------
\ Core storage sizes
\ ----------------------------
10000 CONSTANT MAX-OBJ
20000 CONSTANT MAX-EVT

\ ----------------------------
\ Object allocator
\ ----------------------------
VARIABLE NEXT-ID

\ ----------------------------
\ Union-Find parent + rank
\ ----------------------------
CREATE PARENT MAX-OBJ CELLS ALLOT
CREATE RANK   MAX-OBJ CELLS ALLOT

\ ----------------------------
\ Event log: (tag a b c) per event
\ tag is a small integer
\ ----------------------------
VARIABLE EVT#
CREATE EVT-TAG MAX-EVT CELLS ALLOT
CREATE EVT-A   MAX-EVT CELLS ALLOT
CREATE EVT-B   MAX-EVT CELLS ALLOT
CREATE EVT-C   MAX-EVT CELLS ALLOT
\end{verbatim}

\subsection{Initialization and Allocation}

Initialization clears counters and sets each object to be its own parent.

\begin{verbatim}
: INIT ( -- )
  0 NEXT-ID !
  0 EVT# !
;

: NEW ( -- id )
  NEXT-ID @ DUP 1+ NEXT-ID !
  DUP CELLS PARENT + !      \ parent[id] = id
  0 OVER CELLS RANK + !     \ rank[id] = 0
;
\end{verbatim}

The word \texttt{NEW} implements the creation of a fresh entity. In a strict
event-historical reading, creation itself is an event and should be logged. One
can therefore define \texttt{POP} (below) to include allocation, or define
\texttt{NEW} as a convenience word that calls \texttt{POP} with a distinguished
tag.

\subsection{Event Logging}

We define a helper word that appends a record to the event log.

\begin{verbatim}
: LOG-EVT ( tag a b c -- )
  EVT# @ DUP
  >R
  CELLS EVT-TAG + !
  R@ CELLS EVT-A   + !
  R@ CELLS EVT-B   + !
  R> CELLS EVT-C   + !
  EVT# @ 1+ EVT# !
;
\end{verbatim}

This realizes the fundamental Spherepop thesis that semantics should be sourced
from events. Every operator below calls \texttt{LOG-EVT}.

\subsection{Union-Find Core for \textsc{Merge}}

A minimal \textsc{Merge} can be implemented by union-find. This yields an
efficient derived part-of structure while the authoritative meaning remains in
the log.

\begin{verbatim}
: FIND ( x -- root )
  DUP CELLS PARENT + @ OVER = IF EXIT THEN
  DUP CELLS PARENT + @ RECURSE
  DUP ROT CELLS PARENT + !
;

: UNION ( x y -- )
  FIND SWAP FIND
  2DUP = IF 2DROP EXIT THEN
  \ union by rank
  2DUP CELLS RANK + @ SWAP CELLS RANK + @
  2DUP < IF \ rank[x] < rank[y]
     DROP SWAP \ ensure first has >= rank
  ELSE DROP THEN
  \ now: xRoot yRoot with rank[x]>=rank[y]
  OVER CELLS PARENT + !  \ parent[yRoot] = xRoot
  \ if ranks equal, increment rank
  2DUP CELLS RANK + @ SWAP CELLS RANK + @ = IF
     OVER CELLS RANK + DUP @ 1+ SWAP !
  THEN
  2DROP
;
\end{verbatim}

\subsection{Spherepop Operators as Forth Words}

We now define Spherepop operators. The semantics here are intentionally minimal
and should be read as a reference architecture rather than a completed theory.
The key feature is that every operator both (i) performs a derived-state update
and (ii) logs itself as an event for replay.

\subsubsection{\textsc{Pop}}

A \textsc{Pop} is an irreversible event that introduces a new object and
thereby changes the future space by committing to a new historical fact: the
object exists and can be referred to. One may also interpret \textsc{Pop} as a
commitment that excludes certain future histories in which the object did not
exist.

\begin{verbatim}
1 CONSTANT TAG-POP

: POP ( -- id )
  NEW DUP 0 0 0
  TAG-POP ROT ROT ROT LOG-EVT
;
\end{verbatim}

The particular record fields are unused here; in richer variants, \textsc{Pop}
may attach provenance or a parent context.

\subsubsection{\textsc{Bind}}

A \textsc{Bind} introduces ordering or dependency without removing options. In a
minimal implementation, \textsc{Bind} can be represented as a constraint edge in
a separate structure (e.g.\ a dependency graph), while leaving the union-find
structure untouched. For compactness, we log it and omit derived enforcement;
the essential point is the event-historical record.

\begin{verbatim}
2 CONSTANT TAG-BIND

: BIND ( a b -- )
  \ meaning: a must precede b, or a constrains b
  0 TAG-BIND ROT ROT 0 LOG-EVT
;
\end{verbatim}

A full implementation would maintain a constraint graph and enforce acyclicity
or other admissibility criteria. The present definition demonstrates the
operational encoding: binding is recorded and therefore participates in replay.

\subsubsection{\textsc{Merge}}

A \textsc{Merge} identifies two objects into a common equivalence class. In the
derived state, this is union-find; in the log, it is a recorded event.

\begin{verbatim}
3 CONSTANT TAG-MERGE

: MERGE ( a b -- )
  2DUP UNION
  0 TAG-MERGE ROT ROT 0 LOG-EVT
;
\end{verbatim}

\subsubsection{\textsc{Link}}

A \textsc{Link} introduces a relation between objects without identifying them.
This can be modeled as an adjacency list or as a set of labeled edges. Again we
log the event and supply a minimal derived mechanism.

Here we assume a fixed-size edge store. The representation can be improved, but
the essential semantics are: link is explicit, replayable, and historically
grounded.

\begin{verbatim}
4 CONSTANT TAG-LINK
100000 CONSTANT MAX-LINK
VARIABLE LNK#
CREATE LNK-FROM MAX-LINK CELLS ALLOT
CREATE LNK-TO   MAX-LINK CELLS ALLOT

: LINK-ADD ( a b -- )
  LNK# @ DUP >R
  R@ CELLS LNK-FROM + !
  R> CELLS LNK-TO   + !
  LNK# @ 1+ LNK# !
;

: LINK ( a b -- )
  2DUP LINK-ADD
  0 TAG-LINK ROT ROT 0 LOG-EVT
;
\end{verbatim}

\subsubsection{\textsc{Collapse}}

A \textsc{Collapse} discards historical differentiation while preserving
flexibility. Operationally, collapse is best understood as log compression:
replace a detailed prefix of events with a summary that yields the same derived
state for the purposes of future evolution, while erasing the fine-grained
historical distinctions.

In Forth, the simplest demonstration is to reset the log to a checkpoint while
retaining the derived structures. This does not undo the derived constraints; it
forgets the path by which they were achieved. More sophisticated collapse can
store a digest event that records only invariants.

\begin{verbatim}
5 CONSTANT TAG-COLLAPSE
VARIABLE EVT-CUT

: CHECKPOINT ( -- )
  EVT# @ EVT-CUT !
;

: COLLAPSE ( -- )
  \ Forget events prior to EVT-CUT by truncation.
  \ Derived state is kept; only historical detail is removed.
  EVT-CUT @ EVT# !
  0 0 0 TAG-COLLAPSE 0 0 0 LOG-EVT
;
\end{verbatim}

The call to \texttt{LOG-EVT} after truncation functions as a marker indicating
that a collapse occurred. In a richer system, the collapse record would contain
a hash of the retained derived state or a canonical summary of invariants.

\subsection{Replay and Verification}

Replay is central to Spherepop semantics. In this implementation, replay means
resetting the derived state and reapplying the event log to reconstruct the
present. The following word replays events, interpreting each tag.

\begin{verbatim}
: RESET-DERIVED ( -- )
  \ Reset union-find for all allocated objects.
  NEXT-ID @ 0 DO
    I DUP CELLS PARENT + !
    0 I CELLS RANK + !
  LOOP
  0 LNK# !
;

: EVT@ ( i -- tag a b c )
  DUP CELLS EVT-TAG + @
  OVER CELLS EVT-A   + @
  OVER CELLS EVT-B   + @
  SWAP CELLS EVT-C   + @
;

: APPLY ( tag a b c -- )
  \ interpret a single event
  >R >R >R \ stash c b a
  CASE
    TAG-POP OF \ pop: nothing to do; object already allocated in this toy model
      R> R> R> DROP DROP DROP
    ENDOF
    TAG-BIND OF
      R> R> R> DROP DROP DROP
    ENDOF
    TAG-MERGE OF
      R> R> SWAP R> DROP \ a b c -> a b
      UNION
    ENDOF
    TAG-LINK OF
      R> R> SWAP R> DROP
      LINK-ADD
    ENDOF
    TAG-COLLAPSE OF
      R> R> R> DROP DROP DROP
    ENDOF
  ENDCASE
;

: REPLAY ( -- )
  RESET-DERIVED
  EVT# @ 0 DO
    I EVT@ APPLY
  LOOP
;
\end{verbatim}

This replay mechanism exhibits the intended separation: the log is the semantic
authority and the derived structures are reconstructions. As written, \textsc{Pop}
does not allocate during replay; a strict version would log allocation and
reallocate in replay, or store object counts as part of a checkpoint event. The
point is architectural: there is a canonical history and a derived present.

\subsection{Event-Nodes and Texture-Nodes in Forth Terms}

Within this implementation, event-nodes correspond to log entries whose replay
changes the admissible future space. \textsc{Pop} and \textsc{Merge} are event-like
because they introduce new referents or collapse referents irreversibly in the
future algebra. \textsc{Collapse} is also event-like, though of a different kind:
it changes the space of counterfactual explanation by erasing distinctions.

Texture-nodes correspond to parameters that alter generative behavior without
changing what is admissible in principle. In Forth, these are naturally modeled
as variables that affect the interpretation of subsequent events or the
stochasticity/heuristics by which new events are proposed, while leaving replay
equivalence intact. The most faithful implementation strategy is therefore to
treat texture as fiberwise metadata: recorded ratios, seeds, and generator
parameters that can be changed without altering the set of admissible histories,
so long as they satisfy the reconstructibility and support-preservation
conditions derived earlier.

\subsection{Remarks on Safety and Minimality}

This appendix intentionally avoids dynamic code generation, self-modifying
definitions, or arbitrary memory writes. Forth makes such techniques possible,
but Spherepop’s semantic commitment is stronger when the only authoritative
structure is the event log and its replay. The minimal implementation therefore
prefers explicit records and replayable interpreters, because these preserve the
intended distinction between structural history and derived present state.

In this sense, Forth serves as an existence proof: Spherepop operators can be
realized in a compact, transparent substrate where mechanistic inspection is
straightforward and where the semantics of irreversibility are expressed as
literal log structure rather than as implicit side effects.

% ==================================================
\section{Appendix: Implementing Spherepop Operators in Haskell}
% ==================================================

This appendix provides a Haskell implementation sketch for Spherepop operators
analogous in spirit to the preceding Forth appendix. The guiding design
constraint is that the system’s semantics be \emph{event-sourced}: the event log
is authoritative, while the present state is a derived object obtained by replay
or by replay-equivalent compression. The implementation therefore separates (i)
the type of events, (ii) an interpreter (replay function) from events to state,
and (iii) a small API of Spherepop operators that append to the log.

The code below is intentionally minimal and intended as a reference
architecture. It is designed to be readable, auditable, and compatible with
pure functional reasoning. Performance and storage layout can be improved later
without changing the semantic core.

\subsection{Core Types: Objects, Events, and Histories}

We represent objects as integer identifiers. A history is a list of events. The
event type is explicit and closed; this is the key to keeping semantics
inspectable.

\begin{verbatim}
{-# LANGUAGE DeriveGeneric #-}

module Spherepop where

import GHC.Generics (Generic)
import Data.Map.Strict (Map)
import qualified Data.Map.Strict as M
import Data.Set (Set)
import qualified Data.Set as S

-- | Object identifiers.
newtype Obj = Obj { unObj :: Int }
  deriving (Eq, Ord, Show, Generic)

-- | Spherepop events. The log is authoritative.
data Event
  = Pop Obj
  | Bind Obj Obj
  | Merge Obj Obj
  | Link Obj Obj
  | Collapse Checkpoint
  deriving (Eq, Show, Generic)

-- | A checkpoint can be as small as a log index plus an optional digest.
data Checkpoint = Checkpoint
  { cutIndex :: Int
  , digest   :: Maybe String
  } deriving (Eq, Show, Generic)

-- | A history is a sequence of events.
newtype History = History { unHistory :: [Event] }
  deriving (Eq, Show, Generic)
\end{verbatim}

\subsection{Derived State}

A derived state is an interpretation of the history. We implement a minimal
state consisting of a union-find style representative map for merges, a set of
links, and an optional set of binding constraints. Bind constraints are recorded
as data; enforcing them (e.g.\ checking acyclicity) can be layered on later.

\begin{verbatim}
-- | Representative mapping for Merge semantics.
-- Invariant: rep[x] is the chosen representative for x's class.
type Rep = Map Obj Obj

-- | Link relation as a set of directed edges.
type Links = Set (Obj, Obj)

-- | Bind constraints as a set of directed edges (ordering or dependency).
type Binds = Set (Obj, Obj)

data State = State
  { reps  :: Rep
  , links :: Links
  , binds :: Binds
  , next  :: Int
  } deriving (Eq, Show)

emptyState :: State
emptyState = State
  { reps  = M.empty
  , links = S.empty
  , binds = S.empty
  , next  = 0
  }
\end{verbatim}

The field \texttt{next} provides fresh identifiers. In a strict event-historical
semantics, allocation should itself be recorded as an event. We do so by
ensuring that \texttt{pop} constructs the fresh object and records it explicitly.

\subsection{Union-Find Operations}

For minimality, we implement a simple representative-chasing \texttt{find}
operation with path compression. Union chooses the representative deterministically.

\begin{verbatim}
findRep :: Rep -> Obj -> Obj
findRep r x =
  case M.lookup x r of
    Nothing -> x
    Just p  ->
      let root = findRep r p
      in root

compress :: Rep -> Obj -> Rep
compress r x =
  case M.lookup x r of
    Nothing -> r
    Just p  ->
      let root = findRep r p
      in M.insert x root r

unionRep :: Rep -> Obj -> Obj -> Rep
unionRep r a b =
  let ra = findRep r a
      rb = findRep r b
  in if ra == rb
     then r
     else
       -- deterministically pick smaller Obj as representative
       let rep = min ra rb
           sub = max ra rb
           r'  = M.insert sub rep r
       in compress (compress r' a) b
\end{verbatim}

\subsection{Replay Semantics: Interpreting Events}

Replay is the semantic heart: given a history, compute the derived state. This
function is pure and therefore mechanically checkable. In particular, any two
histories that replay to the same state are equivalent at the chosen level of
derived representation.

\begin{verbatim}
apply :: Event -> State -> State
apply ev st =
  case ev of
    Pop o ->
      -- Ensure the object exists in reps map (optional) and advance allocation.
      st { reps = M.insertWith (\_ old -> old) o o (reps st) }

    Bind a b ->
      st { binds = S.insert (a,b) (binds st) }

    Merge a b ->
      st { reps = unionRep (reps st) a b }

    Link a b ->
      st { links = S.insert (a,b) (links st) }

    Collapse cp ->
      -- Collapse is handled at the history level; apply is a no-op here.
      -- We still record its existence in the log, and replay interprets it
      -- as a truncation marker when requested.
      st

replay :: History -> State
replay (History evs) =
  foldl (flip apply) emptyState evs
\end{verbatim}

The \texttt{Collapse} event is intentionally treated as a semantic instruction
about the history rather than as a state mutation. This mirrors the idea that
collapse erases historical differentiation while preserving the derived present.

\subsection{Collapse as Log Compression}

We implement collapse as truncation to a checkpoint index followed by insertion
of a \texttt{Collapse} marker. More sophisticated implementations can replace
the truncated prefix with a digest event that stores canonical invariants.

\begin{verbatim}
collapse :: Checkpoint -> History -> History
collapse cp (History evs) =
  let k   = max 0 (min (cutIndex cp) (length evs))
      evs' = take k evs
  in History (evs' ++ [Collapse cp])
\end{verbatim}

If one prefers that collapse preserve derived state exactly while discarding
path detail, then the checkpoint should be taken at a point where replayed state
is already acceptable as a canonical summary. Alternatively, one can store a
digest that commits to a particular derived state and treat the digest as an
event-node in the sense of the main text.

\subsection{Spherepop Operators as an API}

We now define a small API that maintains a history and a derived state in a
single structure. In practice, the derived state is optional because it can
always be recomputed from replay, but keeping it avoids repeated folds.

\begin{verbatim}
data World = World
  { hist  :: History
  , state :: State
  } deriving (Eq, Show)

emptyWorld :: World
emptyWorld = World (History []) emptyState

append :: Event -> World -> World
append ev (World (History evs) st) =
  let st' = apply ev st
  in World (History (evs ++ [ev])) st'

-- | POP allocates a new object and records it.
pop :: World -> (Obj, World)
pop (World h st) =
  let o   = Obj (next st)
      st' = st { next = next st + 1 }
      w'  = append (Pop o) (World h st')
  in (o, w')

bind :: Obj -> Obj -> World -> World
bind a b = append (Bind a b)

merge :: Obj -> Obj -> World -> World
merge a b = append (Merge a b)

link :: Obj -> Obj -> World -> World
link a b = append (Link a b)

-- | COLLAPSE compresses history at the log level, then rebuilds derived state.
collapseWorld :: Checkpoint -> World -> World
collapseWorld cp (World h _) =
  let h'  = collapse cp h
      st' = replay h'
  in World h' st'
\end{verbatim}

This API makes explicit what the event-historical semantics demands: all
meaningful changes are events, and the present is derived. The only non-eventful
change above is \texttt{next}, the allocator counter; this is an implementation
detail. If allocation must itself be replayable from events, then the allocator
counter should be computed from the history (e.g.\ by counting \texttt{Pop}
events), and \texttt{State.next} can be removed.

\subsection{Event-Nodes and Texture-Nodes as Types}

The main body of the essay distinguishes event-nodes from texture-nodes. In the
Haskell setting, the distinction becomes a typing discipline: events are logged
and replayable, while textures are parameters used by interpreters that modulate
behavior without changing admissible futures.

A minimal way to represent this is to parameterize the interpreter by a texture
record. Texture does not appear in the history; it appears in the semantics of
how events are interpreted or proposed.

\begin{verbatim}
data Texture = Texture
  { noiseSeed  :: Int
  , ratioParam :: Double
  } deriving (Eq, Show)

-- An interpreter that is allowed to consult texture.
applyWithTexture :: Texture -> Event -> State -> State
applyWithTexture _ ev st =
  apply ev st
\end{verbatim}

In the simplest case, texture is irrelevant to replay, and therefore has no
effect on admissible futures; it merely affects how new events are generated.
In richer systems, texture may modulate stochastic proposals, heuristic
selection of merges, or prioritization of links. The necessary and sufficient
conditions derived earlier then become properties of \texttt{applyWithTexture}:
if varying \texttt{Texture} changes reachability support or introduces hysteresis
not reconstructible at the causal boundary, the texture must be promoted into
the event log as an event-node.

\subsection{A Small Worked Example}

\begin{verbatim}
example :: World
example =
  let (a,w1) = pop emptyWorld
      (b,w2) = pop w1
      w3     = link a b w2
      w4     = merge a b w3
      w5     = bind a b w4
      w6     = collapseWorld (Checkpoint 3 Nothing) w5
  in w6
\end{verbatim}

The example demonstrates the intended separation. Events are appended, replay
reconstructs, and collapse truncates history while preserving a derived present.
The meaning of the system is the history; the state is its current projection.

\subsection{Remarks on Extensibility}

This Haskell design is deliberately conservative. It favors explicit event
types and a single interpreter rather than embedding Spherepop semantics into an
opaque monad. Nevertheless, it scales naturally. One can add typed variants,
invariants, and proof-carrying summaries by refining the \texttt{Event} type and
the replay function. One can also add categorical structure by interpreting
histories as morphisms in a free category generated by events, with \texttt{replay}
providing a functor into a concrete state category. Finally, one can implement
sheaf-like texture structure by treating \texttt{Texture} as a fiberwise
assignment whose gluing conditions correspond to reconstructibility constraints.

In all cases, the central Spherepop principle is preserved: the authoritative
meaning is an event history, and the operational present is a derived, replayable
projection of that history.
\end{verbatim}

% ==================================================
\section{Appendix: Implementing Spherepop Operators in Racket}
% ==================================================

This appendix presents a Racket implementation sketch of Spherepop operators
parallel to the Forth and Haskell appendices. The architectural constraint is
again event-sourcing: the event log is the semantic authority, while any
``current state'' is a derived projection obtained by interpreting (replaying)
that log. Racket is well-suited for this because it provides algebraic data via
structs, persistent maps/sets, and a macro system for building a compact surface
syntax without changing the underlying semantics.

\subsection{Core Representation}

We represent objects as integers. Events are represented as transparent structs
so that they are printable and inspectable. A history is a list of events. The
derived state is a record containing representative mappings (for merge), link
edges, bind constraints, and an allocation counter. As in the other appendices,
the allocation counter may be eliminated if one insists that allocation be
computed solely from the history; it is kept here to keep the operational API
simple.

\begin{verbatim}
#lang racket

(require racket/set
         racket/match)

;; ----------------------------
;; Events
;; ----------------------------

(struct pop      (id)       #:transparent)
(struct bind     (a b)      #:transparent)
(struct merge    (a b)      #:transparent)
(struct link     (a b)      #:transparent)
(struct collapse (cut digest) #:transparent)

;; A history is a list of events.
;; (listof (U pop bind merge link collapse))

;; ----------------------------
;; Derived State
;; ----------------------------

(struct state (reps links binds next) #:transparent)

(define empty-state
  (state (hash) (set) (set) 0))
\end{verbatim}

\subsection{Union-Find Style Representatives}

We implement a representative map for merges. This is a minimal union-find
variant expressed in purely functional style. The representative of an object is
obtained by chasing pointers in a hash map. Path compression can be added, but
the minimal version suffices for semantic clarity.

\begin{verbatim}
(define (find-rep reps x)
  (define p (hash-ref reps x #f))
  (if (or (not p) (= p x))
      x
      (find-rep reps p)))

(define (union-rep reps a b)
  (define ra (find-rep reps a))
  (define rb (find-rep reps b))
  (cond [(= ra rb) reps]
        [else
         (define rep (min ra rb))
         (define sub (max ra rb))
         (hash-set reps sub rep)]))
\end{verbatim}

\subsection{Replay Semantics}

Replay interprets a history as a fold over events. This gives a canonical meaning
to a history: two histories are equivalent at the derived level if replay yields
the same state (or the same invariants of state).

Collapse is treated as a history-level operation rather than a state mutation.
In the replay function we therefore ignore collapse events, except insofar as
they are used as markers in separate compression procedures.

\begin{verbatim}
(define (apply-event ev st)
  (match ev
    [(pop id)
     (define reps (state-reps st))
     (define reps* (if (hash-has-key? reps id)
                       reps
                       (hash-set reps id id)))
     (struct-copy state st [reps reps*])]

    [(bind a b)
     (struct-copy state st [binds (set-add (state-binds st) (cons a b))])]

    [(merge a b)
     (struct-copy state st [reps (union-rep (state-reps st) a b)])]

    [(link a b)
     (struct-copy state st [links (set-add (state-links st) (cons a b))])]

    [(collapse _ _)
     st]))

(define (replay hist)
  (for/fold ([st empty-state])
            ([ev (in-list hist)])
    (apply-event ev st)))
\end{verbatim}

\subsection{History-Level Collapse (Compression)}

Collapse is defined as erasing historical differentiation while preserving the
derived present. Operationally, the minimal implementation truncates the log to a
cut index and appends a collapse marker. A more faithful implementation would
replace the truncated prefix with a digest event encoding canonical invariants.
The digest field is included for that purpose, even if unused here.

\begin{verbatim}
(define (collapse-history hist cut [digest #f])
  (define cut* (max 0 (min cut (length hist))))
  (append (take hist cut*)
          (list (collapse cut* digest))))
\end{verbatim}

\subsection{World API: Event-Sourcing with Derived Cache}

As in the other appendices, we provide a \texttt{world} record containing both
history and a cached derived state. Every operator appends an event and updates
the cache by interpreting that single event. Collapse rebuilds the cache by
replaying the compressed history, reflecting the idea that collapse is a
history transformation rather than an ordinary state transition.

\begin{verbatim}
(struct world (hist st) #:transparent)

(define empty-world
  (world '() empty-state))

(define (append-event w ev)
  (define hist* (append (world-hist w) (list ev)))
  (define st*   (apply-event ev (world-st w)))
  (world hist* st*))

(define (POP w)
  (define id (state-next (world-st w)))
  (define st* (struct-copy state (world-st w) [next (add1 id)]))
  (define w*  (world (world-hist w) st*))
  (values id (append-event w* (pop id))))

(define (BIND w a b)
  (append-event w (bind a b)))

(define (MERGE w a b)
  (append-event w (merge a b)))

(define (LINK w a b)
  (append-event w (link a b)))

(define (COLLAPSE w cut [digest #f])
  (define hist* (collapse-history (world-hist w) cut digest))
  (world hist* (replay hist*)))
\end{verbatim}

\subsection{Texture Parameters as Fiberwise Semantics}

The main essay distinguishes event-nodes from texture-nodes. In this Racket
realization, the minimal way to express that distinction is to keep textures
out of the history and instead treat them as parameters of interpretation or of
event proposal. That is, textures are not logged events; they are fiberwise
metadata that modulate generative behavior without changing admissible futures,
unless they violate the support-preservation and reconstructibility conditions
derived in the main text.

One can express this by extending the interpreter to consult a texture record
when applying events, or by making event-proposal functions depend on textures.

\begin{verbatim}
(struct texture (seed ratio) #:transparent)

(define default-texture
  (texture 0 1.0))

(define (apply-event/texture tex ev st)
  ;; In the simplest form, texture does not affect replay semantics at all.
  ;; More elaborate versions can allow texture to bias heuristic choices,
  ;; while remaining fiberwise if it preserves future support.
  (apply-event ev st))
\end{verbatim}

If texture begins to gate regime changes, introduce hysteresis not reconstructible
from the causal boundary, or change the reachable support, then it can no longer
be treated as a mere parameter; it must be promoted into the event log as an
event-node, for instance by introducing a new event constructor such as
\texttt{(set-texture tex)} and replaying it as part of the authoritative history.

\subsection{A Small Worked Example}

\begin{verbatim}
(define (example)
  (define-values (a w1) (POP empty-world))
  (define-values (b w2) (POP w1))
  (define w3 (LINK w2 a b))
  (define w4 (MERGE w3 a b))
  (define w5 (BIND w4 a b))
  (define w6 (COLLAPSE w5 3))
  w6)

(example)
\end{verbatim}

The example illustrates the intended semantics. All meaning-bearing actions are
events appended to a log. Replay reconstructs the present. Collapse compresses
history while preserving a derived state. The design is deliberately explicit:
every important object is a plain struct, and every semantic transformation is a
total function over histories.

\subsection{Macro Surface Syntax (Optional)}

Racket permits a compact surface syntax without changing the semantic model. One
can introduce syntactic forms that expand into calls to \texttt{POP}, \texttt{BIND},
\texttt{MERGE}, \texttt{LINK}, and \texttt{COLLAPSE}. This provides an ergonomic
``Spherepop script'' language while retaining the same replay semantics. The
macro layer is intentionally omitted here to keep the appendix focused on the
semantic core, but the presence of transparent event structs makes such a layer
straightforward: a macro expands into event-appending calls, and the log remains
the canonical source of truth.

In this sense, the Racket implementation serves the same purpose as the Forth
and Haskell sketches: it demonstrates that Spherepop operators are not merely
philosophical primitives but can be realized as a small, inspectable, replayable
event calculus in a conventional programming environment.
\end{verbatim}



\end{document}
